\documentclass[11pt]{article}

\usepackage[a4paper,margin=29mm]{geometry}
\usepackage[T1]{fontenc}
\usepackage{lmodern}
\usepackage{microtype}
\usepackage{amsmath,amssymb,amsthm,amscd,mathtools,mathrsfs}
\usepackage{booktabs}
\usepackage{multicol}
\usepackage[hidelinks]{hyperref}

\numberwithin{equation}{section}
\newtheorem{theorem}{Theorem}[section]
\newtheorem{proposition}[theorem]{Proposition}
\newtheorem{lemma}[theorem]{Lemma}
\newtheorem{corollary}[theorem]{Corollary}
\theoremstyle{definition}
\newtheorem{definition}[theorem]{Definition}
\theoremstyle{remark}
\newtheorem{remark}[theorem]{Remark}

\newcommand{\A}{\mathbf A}
\newcommand{\C}{\mathbf C}
\newcommand{\F}{\mathbf F}
\newcommand{\Q}{\mathbf Q}
\newcommand{\Z}{\mathbf Z}
\newcommand{\PP}{\mathbf P}
\newcommand{\OO}{\mathcal O}
\newcommand{\disc}{\operatorname{disc}}
\newcommand{\Or}{\operatorname{Or}}
\newcommand{\Gal}{\operatorname{Gal}}

\title{The Discriminant Resolvent of the \(A_5\)-Cubic Pencil}
\author{Tavis Rudd}
\date{August 2026}

\begin{document}
\maketitle

\begin{abstract}
The nonstandard \(A_5\)-invariant pencil of cubic threefolds carries five
principal gluing kernels: three rational and two golden.  We prove that this
is not an accidental \(3+2\) count.  The five kernels are the two proper
transitive resolvents of the \(S_3\)-torsor of ordered nonzero two-torsion on
the relative elliptic norm axis.  The rational triple is the root cover and
the golden pair has the sign character, hence is the discriminant double
cover of the same two-division cubic.

Writing \(t\) for the signed cubic parameter and \(T=81t^2\), the unmarked
elliptic period is the standard \(X_0(3)\)-coordinate
\[
 j=\frac{(T+27)(T+3)^3}{T}.
\]
The signed cubic line is the sign modular curve \(r^2=T\), with \(r=9t\)
after choosing its deck convention.  The degree-three resolvent is
\(X_0(6)\), and their common degree-six Galois closure belongs to
\(\Gamma_0(3)\cap\Gamma(2)\).  We give rational equations for all three
covers and distinguish the two modular cusps from the additional interior
boundary points where the cubic degenerates.  In particular the modular
identification concerns the actual relative principal-kernel packet, not
only an auxiliary elliptic curve with the same \(j\)-invariant.  Over the
signed line, the resulting cyclic cubic cover is totally ramified exactly at
the two fibres with five \(A_2\)-points.  At the chordal value \(h=5\), the
hyperelliptic limit has Jacobian isogenous to the fifth power of the elliptic
curve selected by the same norm-axis coordinate.
\end{abstract}

\noindent\textbf{Keywords.} cubic threefold; intermediate Jacobian;
elliptic two-torsion; discriminant resolvent; modular curve; \(A_5\).

\section{The result}

Let \(W\subset\C^6\) be the hyperplane \(\sum_{i=0}^5z_i=0\).  The
nonstandard transitive copy of \(A_5\subset S_6\) has two orbits
\(\mathcal T_+\) and \(\mathcal T_-\), each of size ten, on the three-element
subsets of \(\{0,\ldots,5\}\).  Put
\[
 P=\sum_{i=0}^5z_i^3,
 \qquad
 Q=\sum_{I\in\mathcal T_+}\prod_{i\in I}z_i
       -\sum_{I\in\mathcal T_-}\prod_{i\in I}z_i,
 \qquad
 X_t=\{P+tQ=0\}\subset\PP(W).
\]
The outer coset of the normalizer \(S_5=N_{S_6}(A_5)\) fixes \(P\), negates
\(Q\), and therefore sends \(X_t\) to \(X_{-t}\).  Thus \(t\) is the signed
parameter and \(t^2\) is the unmarked parameter.

On the smooth marked base, the six distinguished dihedral axes in the
intermediate Jacobian form one relative elliptic norm axis \(\mathcal E\).
The integral six-axis construction of \cite[Relative six-axis source and
Principal gluing packet]{RuddEpilogue} identifies the five principal kernels
with the finite local system
\begin{equation}\label{eq:packet}
 \mathscr P=PP_{\F_4}(\F_4\otimes_{\F_2}\mathcal E[2]).
\end{equation}
The comparison in Proposition~\ref{prop:actual-axis} below records why the
elliptic curve used to compute the period is this actual norm axis at the
level needed in \eqref{eq:packet}.

\begin{theorem}[Cubic--modular resolvent]\label{thm:main}
Set \(T=81t^2\).  Then the following assertions hold over the smooth cubic
base.
\begin{enumerate}
\item The elliptic norm axis has
\[
 j(\mathcal E)=J(T):=\frac{(T+27)(T+3)^3}{T}.
\]
Consequently the unmarked period map identifies its normalized coarse base
with an open subcurve of \(X_0(3)\).
\item The packet \(\mathscr P\) is the disjoint union of the root cover
\(\PP(\mathcal E[2])\), of degree three, and a degree-two cover having the
sign character of \(\Gal(\mathcal E[2])\subset S_3\).  The second cover is
the discriminant resolvent of the two-division cubic.
\item Over \(X_0(3)\), the root cover is \(X_0(6)\), the sign cover has
function field \(\Q(T)(\sqrt T)\), and their common Galois closure belongs to
\(\Gamma_0(3)\cap\Gamma(2)\).  After one choice of deck involution, the
signed cubic coordinate identifies with the sign coordinate by
\[
 r^2=T,\qquad r=9t.
\]
\item The compact sign curve has genus zero, two cusps of widths \(2\) and
\(6\), and two elliptic points of order three.  Its map to \(X_0(3)\) is
ramified at the two cusps and nowhere else as a coarse map.
\item The geometric mod-two monodromy of the actual norm axis over the marked
smooth cubic base is exactly \(A_3\simeq C_3\).
\item Over the signed cubic line, the cyclic cubic splitting cover is totally
ramified exactly at the two points \(t^2=-1/3\) and is unramified at every
other cubic boundary point.
\end{enumerate}
The smooth cubic base is obtained from the sign modular curve by deleting,
in addition to the cusps over \(T=0,\infty\), the points over
\(T=-27\) and \(T=729/5\).  Those last two values are interior modular
points, not cusps.
\end{theorem}

The theorem identifies local systems and their monodromy.  It does not say
that a fivefold gluing is literally an elliptic \(2\)-isogeny, nor that the
compactified cubic pencil is the modular stack.

\section{The finite \texorpdfstring{\(3+2\)}{3+2} resolvent}

Let \(V\) be a two-dimensional \(\F_2\)-vector space and
\(V_4=\F_4\otimes_{\F_2}V\).  Scalar extension gives a distinguished
three-element subset
\[
 \PP_{\F_2}(V)\subset\PP_{\F_4}(V_4).
\]
Write \(D(V)\) for its two-element complement.

\begin{lemma}[Sign lemma]\label{lem:sign}
Under \(\operatorname{GL}(V)\simeq S_3\), the action on \(D(V)\) is the sign
of the natural action on \(\PP(V)\).  Thus the two finite \(S_3\)-sets are
\[
 \PP(V)\simeq S_3/C_2,
 \qquad
 D(V)\simeq S_3/A_3.
\]
\end{lemma}

\begin{proof}
Choose a basis and write
\[
 \PP^1(\F_4)=\{\infty,0,1,\omega,\omega^2\},\qquad
 \PP^1(\F_2)=\{\infty,0,1\}.
\]
The transformations \(x\mapsto x+1\) and
\(x\mapsto 1/(x+1)\) generate \(\operatorname{PGL}_2(\F_2)\).  The first
acts as \((0\ 1)(\omega\ \omega^2)\); the second cycles
\((\infty\ 0\ 1)\) and fixes \(\omega,\omega^2\).  The first generator is
odd on each of the three-set and the two-set, while the second is even on
both.  This proves the assertion.
\end{proof}

For a degree-three finite \(\acute etale\) cover \(C\to S\), let
\(\Or(C/S)\) denote the double cover associated to the sign of its monodromy
in \(S_3\).  Lemma~\ref{lem:sign} proves
\begin{equation}\label{eq:torsor-class}
 [D(V)]=[\Or(\PP(V)/S)]\quad\text{in }H^1(S,C_2).
\end{equation}
This equality of torsor classes is canonical.  A sheetwise isomorphism of
the two covers has two choices, interchanged by the deck involution.  This
minor distinction is what makes the normalization \(r=9t\), rather than
\(r=-9t\), a convention.

There is a useful higher-prime form of this finite geometry.  If \(V\) has
dimension two over \(\F_\ell\), then, as \(\operatorname{PGL}(V)\)-sets,
\[
 \PP^1(\F_{\ell^2})
 =\PP^1(\F_\ell)\sqcup
   \bigl(\PP^1(\F_{\ell^2})\setminus\PP^1(\F_\ell)\bigr)
 \simeq \operatorname{PGL}(V)/B
       \sqcup\operatorname{PGL}(V)/C_{\rm ns},
\]
where \(B\) is a Borel subgroup and \(C_{\rm ns}\) a nonsplit Cartan.
Frobenius pairs the second orbit, and its quotient is
\(\operatorname{PGL}(V)/N_{\rm ns}\), where \(N_{\rm ns}\) is the Cartan
normalizer;
compare the modular interpretation in
\cite[Section~2]{RebolledoWuthrich}.  For
\(\ell=2\), the normalizer is all of
\(\operatorname{PGL}_2(\F_2)\), while
\(C_{\rm ns}=A_3\).  Thus only at two does the entire nonsplit-Cartan orbit
reduce to a double cover, namely the sign torsor above.

Now let an elliptic curve be given locally, with \(2\) invertible, by
\[
 E:y^2=x^3+Ax+B.
\]
The nonzero two-torsion points are the roots of
\(f_2(x)=x^3+Ax+B\).  The orientation cover of this root cover is therefore
\[
 u^2=\disc(f_2),
 \qquad
 \disc(f_2)=-4A^3-27B^2.
\]
Since \(\Delta_E=16\disc(f_2)\), it has the same square class as the elliptic
discriminant.  Invariantly, \(\Delta_E\) is a nowhere-vanishing section of
the twelfth power of the Hodge line; the double cover is the torsor of its
square roots.  No global scalar equation for \(\Delta_E\) is implied before
a Hodge trivialization is chosen.

Applying the lemma to \eqref{eq:packet} proves Theorem~\ref{thm:main}(2).
It also gives a useful monodromy consequence: choosing one golden sheet
reduces the mod-two image from \(S_3\) to
\(A_3\simeq C_3\).

\section{The cubic parameter and the actual norm axis}

The period calculation starts with the Prym quotient used by van Geemen and
Yamauchi.  In their coordinate \(u\), two auxiliary parameters in the
Weierstrass model of its elliptic factor are
\[
 a(u)=-\frac{32(u-3)^4}{9(u+1)^3(u+3)^2},\qquad
 b(u)=-\frac{8(u-3)^2(u-1)}{(u+1)(u+3)^2}.
\]
In the notation of \cite[Proposition 3.2]{VGY}, the generalized Weierstrass
coefficients are
\begin{align*}
 a_1&=b^2/4+b-4,&
 a_3&=(a+b)^2/2+4a+6b-8,\\
 a_2&=2-a-ab/2-b^4/64-b^3/8-b^2/4-b,&
 a_4&=4(a+b-1),\\
 a_6&=(a+b-1)(8-4a-2ab-b^4/16-b^3/2-b^2-4b).
\end{align*}
The Fourier change of coordinates from the six-point equation gives
\(u=ct\), where \(c^2=5\).  Substitution in the standard formulas for
\(c_4\) and \(\Delta\), followed by \(u^2=5t^2\), yields
\begin{equation}\label{eq:j-t}
 j=\frac{9(3t^2+1)(27t^2+1)^3}{t^2}.
\end{equation}
This calculation is replayed exactly in the paper's verification bundle.

Putting \(T=81t^2\), equation \eqref{eq:j-t} becomes
\begin{equation}\label{eq:J}
 j=J(T)=\frac{(T+27)(T+3)^3}{T}.
\end{equation}
This is the standard degree-four \(j\)-map on \(X_0(3)\).  A convenient
Tate normal form is
\begin{equation}\label{eq:tate}
 E_T:\quad y^2+(T+27)xy+(T+27)^2y=x^3.
\end{equation}
The point \((0,0)\) has order three, and the standard Weierstrass formulas
give
\begin{equation}\label{eq:tate-invariants}
 \Delta(E_T)=T(T+27)^8,\qquad j(E_T)=J(T).
\end{equation}

The equality of \(j\)-maps alone would identify only a generic elliptic
isomorphism class.  The next proposition supplies the geometric comparison
needed for the principal-kernel packet.

\begin{proposition}[Actual-axis comparison]\label{prop:actual-axis}
Over the common marked smooth \(A_5/D_5\) base, the van Geemen--Yamauchi
Prym factor maps isomorphically, as a polarized
elliptic scheme, to the
primitive dihedral norm axis inside the intermediate Jacobian.  In
particular its two-torsion local system is the \(\mathcal E[2]\) occurring in
\eqref{eq:packet}.  Relative to \eqref{eq:tate}, it may differ by the
quadratic twist
\[
 D(T)=(T+27)(T-729/5),
\]
which does not alter its two-torsion.
\end{proposition}

\begin{proof}
The degree-five cyclic quotient construction of
\cite[Propositions 2.1, 3.1, and 3.2]{VGY} gives the Prym elliptic curve and
a pullback map to the corresponding norm-fixed elliptic subvariety of the
intermediate Jacobian.  Pullback followed by norm is multiplication by five,
so the image is the primitive one-dimensional fixed factor.  On the Prym,
the principal polarization pulls back as five times its intrinsic principal
polarization.  The restriction of the intermediate-Jacobian polarization to
the primitive norm axis is likewise five times its primitive principal
polarization.  Factoring pullback through that axis and comparing degrees
therefore gives degree one.  This proves the polarized isomorphism.

The explicit change to \eqref{eq:tate} gives the displayed
twist; alternatively, comparing their \(c_4,c_6\) invariants gives the same
square class.  In characteristic different from two, a quadratic twist
changes the Galois representation by a scalar \(\{\pm1\}\); that scalar is
trivial on two-torsion and preserves every cyclic order-three subgroup.
Hence either comparison identifies the required two-division local system,
and the cyclic level-three marking descends even when a chosen generator does
not.
\end{proof}

The proof also explains the role of earlier results.  Roulleau exhibits the
nonstandard \(A_5\)-pencil and its elliptic-power decomposition
\cite[Theorem 11(D) and Lemma 17]{Roulleau}; Hartlieb places its intermediate
Jacobians in a one-dimensional special subvariety
\cite[Lemma 5.5 and Proposition 5.7]{Hartlieb}.  Neither statement by itself
identifies the relative two-division packet used here.

\section{The modular \texorpdfstring{\(S_3\)}{S3}-diagram}

Reduction modulo two gives
\[
 \rho_2:\Gamma_0(3)\longrightarrow
 \operatorname{SL}_2(\F_2)\simeq S_3.
\]
It is surjective: the elements
\[
 \begin{pmatrix}1&1\\0&1\end{pmatrix},\qquad
 \begin{pmatrix}1&0\\3&1\end{pmatrix}
\]
belong to \(\Gamma_0(3)\) and reduce to two distinct transvections, which
generate \(\operatorname{SL}_2(\F_2)\).  Define
\[
 \Gamma_{\mathrm{sgn}}(6)=
 \ker(\operatorname{sgn}\circ\rho_2),
 \qquad
 \Gamma(2,3)=\Gamma(2)\cap\Gamma_0(3).
\]
The first has index two and the second index six in \(\Gamma_0(3)\).
The resulting normalized Cartesian square is
\begin{equation}\label{eq:resolvent-square}
\begin{CD}
 X(\Gamma(2,3)) @>{2}>> X_0(6)\\
 @V{3}VV                    @VV{3}V\\
 X_{\mathrm{sgn}}(6) @>{2}>> X_0(3).
\end{CD}
\end{equation}

A nonzero point of \(E[2]\) is the same as a cyclic subgroup of order two.
Together with the given cyclic subgroup of order three, it determines a
unique cyclic subgroup of order six.  Conversely, an order-six cyclic
subgroup has unique subgroups of orders two and three.  Therefore the
degree-three root resolvent is
\[
 X_0(6)\longrightarrow X_0(3).
\]
The sign resolvent belongs to \(\Gamma_{\mathrm{sgn}}(6)\), and the full
ordered two-torsion cover belongs to \(\Gamma(2,3)\).  This proves the group
part of Theorem~\ref{thm:main}(3).

All three covers are rational.  With \(y\) on the root cover, \(r\) on the
sign cover, and \(v\) on the common splitting cover, one choice of equations
is
\begin{align}
 T&=-\frac{(4y+3)(y+3)^2}{(y+1)^2},\label{eq:root}\\
 x&=-\frac{4y^4}{(y+1)^2},\label{eq:chosen-root}\\
 r^2&=T,\label{eq:sign}\\
 y&=-\frac{v^2+3}{4},&
 T&=\frac{v^2(9-v^2)^2}{(1-v^2)^2},&
 r&=\frac{v(9-v^2)}{1-v^2}.\label{eq:split}
\end{align}
Substitution of \eqref{eq:root} and \eqref{eq:chosen-root} in the
two-division cubic below gives zero, so \(y\) really selects a nonzero
two-torsion point.  The last line directly verifies the root and sign cover
equations.  It realizes the
classical fact that adjoining a root of an \(S_3\)-cubic and adjoining the
square root of its discriminant generate its splitting field.

\begin{corollary}[Golden orientation cyclicizes the rational packet]
After base change to the sign curve, the root cover \(X_0(6)\to X_0(3)\)
becomes a connected Galois cover of degree three with group \(A_3\simeq C_3\).
It is the full splitting cover of the two-division cubic.  Hence, over the
signed smooth cubic base, the rational three-sheet gluing packet has cyclic
monodromy and any one rational sheet determines the other two.
\end{corollary}

\begin{proof}
The sign subgroup has mod-two image \(A_3\), while a point stabilizer in
\(S_3\) has order two.  Their intersection is trivial.  Therefore the
pullback of the point-stabilizer cover to the sign cover has regular
\(A_3\)-action on its three sheets.  Equivalently, the fibre product of
\eqref{eq:root} and \eqref{eq:sign} is the splitting field parametrized by
\eqref{eq:split}.
\end{proof}

\begin{corollary}[Kummer normal form and local monodromy]
Let \(a^2=-3\).  For the coordinates in \eqref{eq:split},
\begin{equation}\label{eq:kummer-cubic}
 r^2+27=\frac{(v^2+3)^3}{(v^2-1)^2},
 \qquad
 \frac{r-3a}{r+3a}=\left(\frac{v-a}{v+a}\right)^3.
\end{equation}
Thus the cyclic cubic cover of the signed line is totally ramified at
\(v=\pm a\), over \(r=\pm3a\), and nowhere else.  These are precisely the
two points \(t=\pm a/3\), where \(t^2=-1/3\) and the cubic has five
\(A_2\)-points.  The two local monodromies are inverse three-cycles; the
mod-two monodromy around the two cusps and the two chordal points is trivial.
The deck group is already rational over \(\Q(r)\): it is generated by
\begin{equation}\label{eq:deck-sigma}
 \sigma(v)=\frac{v-3}{v+1},
 \qquad
 \sigma^2(v)=-\frac{v+3}{v-1},
 \qquad
 \sigma^3(v)=v.
\end{equation}
Reversing the golden orientation sends \((r,v)\) to \((-r,-v)\) and
conjugates \(\sigma\) to \(\sigma^{-1}\).
\end{corollary}

\begin{proof}
Both identities follow by substitution in \eqref{eq:split}.  Equivalently,
\[
 \frac{dr}{dv}=\frac{(v^2+3)^2}{(v^2-1)^2},
\]
so \(v=\pm a\) are the only critical points and each has ramification index
three.  Their branch values are \(r=\pm3a\).  The product relation on the
punctured signed line makes the two three-cycles inverse.  At the cusps, the
split and sign curves have equal widths \(2\) or \(6\), so the map is
unramified; it is also unramified at the ordinary chordal points.  This proves
Theorem~\ref{thm:main}(6).  Direct substitution gives
\(r(\sigma(v))=r(v)\), and the fixed-point equation for \(\sigma\) is
\(v^2=-3\).  Finally
\(-\sigma(-v)=\sigma^{-1}(v)\), which proves the orientation assertion.
\end{proof}

For completeness, the two-division cubic of \eqref{eq:tate}, obtained by
eliminating \(y\) from the equation and
\(2y+(T+27)x+(T+27)^2=0\), is
\[
 4x^3+(T+27)^2x^2+2(T+27)^3x+(T+27)^4.
\]
Its discriminant is
\begin{equation}\label{eq:disc}
 16T(T+27)^8.
\end{equation}
Thus its discriminant square class is exactly \(T\).  By
Proposition~\ref{prop:actual-axis} and Lemma~\ref{lem:sign}, the exotic pair
in the actual gluing packet is \eqref{eq:sign}.  Pulling back along
\(T=81t^2\) splits it; choosing one sheet gives \(r=9t\).  This completes
the proof of Theorem~\ref{thm:main}(1)--(3).  Over the sign curve the mod-two
image is \(A_3\), by construction of \(\Gamma_{\mathrm{sgn}}(6)\).  Passing
to the smaller smooth cubic open cannot decrease that image: concretely, a
small loop around either deleted lift of the order-three point \(T=-27\)
acts as a three-cycle.  Proposition~\ref{prop:actual-axis} transports this
action to the actual norm axis.  This proves assertion~(5).

\section{Compactification and the cubic boundary}

The coarse orbifold \(X_0(3)\) has two cusps of widths one and three and one
elliptic point of order three.  A primitive parabolic at either cusp reduces
modulo two to a transposition, so it does not lie in
\(\Gamma_{\mathrm{sgn}}(6)\), while its square does.  Each cusp consequently
has one ramified lift, of width two or six.  An order-three elliptic generator
reduces to a three-cycle and lies in the sign kernel; the elliptic point has
two unramified lifts, both of order three.

The sign subgroup has index eight in \(\operatorname{PSL}_2(\Z)\).  The
genus formula now reads
\[
 g=1+\frac{8}{12}-\frac{0}{4}-\frac{2}{3}-\frac{2}{2}=0.
\]
This proves Theorem~\ref{thm:main}(4), including the absence of any further
coarse branching.

The same inertia calculation completes the passport of
\eqref{eq:resolvent-square}.  On the root cover, a transposition has orbits of
sizes one and two, giving cusp widths \(1,2,3,6\).  On the regular split
cover it has three orbits of size two, giving widths
\(2,2,2,6,6,6\).  The root and split subgroups contain no elliptic torsion.
Their indices in \(\operatorname{PSL}_2(\Z)\) are twelve and twenty-four, so
the genus formula gives genus zero in both cases.  Thus every curve in the
resolvent square is rational, now independently of the displayed
parametrizations.

The eta Hauptmodul
\[
 h(\tau)=\left(\frac{\eta(\tau)}{\eta(3\tau)}\right)^{12}
\]
has a simple pole at the width-one cusp and a simple zero at the width-three
cusp.  In the normalization \eqref{eq:J},
\[
 h=\frac{729}{T}.
\]
Because \(729=27^2\), the quadratic fields
\(\Q(h)(\sqrt h)\) and \(\Q(T)(\sqrt T)\) coincide.  Analytically one may
take \(\sqrt h=(\eta(\tau)/\eta(3\tau))^6\).  Its half-integral order in the
base cusp parameter is the width doubling just computed.

Combining this square root with \(r=9t\) gives an explicit modular
uniformization of the signed cubic parameter:
\begin{equation}\label{eq:eta-t}
 \boxed{\displaystyle
 t(\tau)=3\left(\frac{\eta(3\tau)}{\eta(\tau)}\right)^6}
\end{equation}
after the same deck convention.  In the \(h\)-coordinate, the two
noncuspidal cubic boundary values below are especially simple:
\[
 T=-27\Longleftrightarrow h=-27,
 \qquad
 T=729/5\Longleftrightarrow h=5.
\]

It remains to compare compactifications.  A direct singularity calculation
for the six-variable cubic equation gives
\begin{center}
\begin{tabular}{lll}
\toprule
cubic value & \(T\)-value & degeneration \\
\midrule
\(t=\infty\) & \(T=\infty\) & six singular points \\
\(t=0\) & \(T=0\) & ten singular points \\
\(t^2=-1/3\) & \(T=-27\) & five \(A_2\)-points \\
\(t^2=9/5\) & \(T=729/5\) & chordal cubic \\
\bottomrule
\end{tabular}
\end{center}

\begin{proposition}[The chordal elliptic factor]\label{prop:chordal-factor}
At \(T=729/5\), the actual transverse term of the pencil cuts a reduced
degree-twelve divisor on the singular rational normal curve of the chordal
cubic.  The limit intermediate Jacobian is the Jacobian of the
icosahedral hyperelliptic curve
\[
 C_{\rm ico}:\quad w^2=x(x^{10}+11x^5-1).
\]
Moreover
\[
 J(C_{\rm ico})\sim E_{\rm ico}^{\,5},
 \qquad
 E_{\rm ico}:\quad y^2=x(x^2+11x-1),
\]
and the elliptic factor has
\[
 j(E_{\rm ico})=\frac{2^{14}31^3}{5^3}
 =J\left(\frac{729}{5}\right).
\]
Consequently the modular value \(h=5\) records the elliptic isogeny factor
of the hyperelliptic chordal limit.
\end{proposition}

\begin{proof}
It is enough to treat one of the two marked points above \(T=729/5\), since
the outer normalizer exchanges them.  Let \(\zeta\) be a primitive fifth root
of unity and put
\[
 p=[1:\zeta:\zeta^2:\zeta^3:\zeta^4:0]\in\PP(W).
\]
For
\[
 t_0=-\frac35-\frac65(\zeta^2+\zeta^3),
 \qquad t_0^2=\frac95,
\]
direct differentiation gives \(p\in\operatorname{Sing}(X_{t_0})\) and \(Q(p)=0\).
Its projective \(A_5\)-orbit \(Z\) has twelve points, all lying on the
singular rational normal quartic
\(C=\operatorname{Sing}(X_{t_0})\), and \(Q\) vanishes on
all of them.

The exact linear-algebra calculation is short.  Evaluation of the fifteen
quadrics on \(Z\) has rank nine, so \(I_Z(2)\) has dimension six.  Since a
rational normal quartic also has \(\dim I_C(2)=15-9=6\), the inclusion
\(I_C(2)\subset I_Z(2)\) is equality.  The ideal of a rational normal curve
is generated by its quadrics.  In degree three their linear multiples span a
space of dimension twenty-two, whereas adjoining \(Q\) raises the rank to
twenty-three.  Thus \(Q\notin I_C(3)\).

Consequently \(Q|_C\) is a nonzero section of
\(\mathcal O_C(3)\simeq\mathcal O_{\PP^1}(12)\).  Its zero divisor contains
the twelve distinct points of \(Z\), so it is exactly \(Z\) and is reduced.
The first-order deformation of \(P+t_0Q\) in this pencil is \(Q\); hence the
nonvanishing criterion of Allcock--Carlson--Toledo shows that this direction
is transverse to the chordal locus \cite[Lemma~2.5]{ACT}.  The general
chordal-limit description therefore applies
\cite[Remark~1.8]{CMGHL}.  The icosahedral action on \(C\simeq\PP^1\) has a
unique orbit of length twelve.  In a standard coordinate it is the zero
divisor of
\[
 su(s^{10}+11s^5u^5-u^{10}),
\]
which gives the displayed equation for \(C_{\rm ico}\).

Paulhus proves the displayed isogeny decomposition and elliptic factor
\cite[Theorem~2]{Paulhus}.  For the printed Weierstrass equation,
\(c_4=1984=2^6\cdot31\) and \(\Delta=2000=2^4\cdot5^3\), hence
\(j(E_{\rm ico})=2^{14}31^3/5^3\).  Direct substitution in
\(J(T)=(T+27)(T+3)^3/T\) gives the same value at \(T=729/5\).
\end{proof}

\begin{corollary}[Boundary reconstruction]\label{cor:boundary-reconstruction}
The embedded branch orbit \(Z\subset\PP(W)\) determines the pair
\((C,[Q|_C])\), hence the chordal central fibre and its projective
first-order normal direction.
\end{corollary}

\begin{proof}
Quadratic generation and \(I_Z(2)=I_C(2)\) recover \(C\) as a scheme; its
secant variety is the chordal cubic.  The divisor
\(Z\) uniquely determines \([Q|_C]\in\PP H^0(C,\mathcal O_C(3))\).
Projective normality identifies this with the normal direction
\[
 [Q]\in\PP\bigl(H^0(\PP(W),\mathcal O(3))/I_C(3)\bigr).
\]
\end{proof}

Only \(T=0\) and \(T=\infty\) are cusps of \(X_0(3)\).  The marked smooth
cubic base is therefore
\[
 \PP^1_t\setminus
 \{0,\infty,\ t^2=-1/3,\ t^2=9/5\},
\]
and its unmarked image is
\[
 X_0(3)\setminus\{0,\infty,-27,729/5\}.
\]
On these opens, \(t\mapsto T=81t^2\) is finite \(\acute etale\).  Its
compactification ramifies only over the two modular cusps.  This proves the
boundary assertion of Theorem~\ref{thm:main} and prevents a misleading
identification of every cubic degeneration with a cusp.

\section{Conclusion and scope}

The object behind the five principal kernels is the \(S_3\)-torsor of ordered
nonzero two-torsion on the actual relative norm axis.  Its two proper
transitive quotient sets are
\[
 S_3/C_2\quad(3\text{ sheets}),
 \qquad
 S_3/A_3\quad(2\text{ sheets}).
\]
Thus the golden pair contains no independent monodromy datum: it is the
discriminant orientation of the rational root triple.  The familiar
decomposition
\(\PP^1(\F_4)=\PP^1(\F_2)\sqcup\{\omega,\omega^2\}\) is the fibrewise
form of this global resolvent theorem.

The modular identification here is deliberately limited.  It neither turns
the five principal kernels into ordinary elliptic isogenies nor promotes the
whole compactified cubic pencil to a modular stack.  The additional cubic
boundary points at \(T=-27\) and \(729/5\) are not modular cusps.  At the
second point Proposition~\ref{prop:chordal-factor} identifies the
hyperelliptic limit and its elliptic factor.  Related level-three period
maps occur for the Winger
pencil \cite[Theorem 1.1 and Proposition 5.1]{LZ}; that is a distinct family and
is not used to identify the cubic base.  The contribution here is the
combination of the explicit cubic period coordinate, the polarized
actual-axis comparison, and the principal-kernel local system into one
resolvent diagram.

\paragraph{Reproducibility.}
The directory \texttt{verification/} contains a deterministic SymPy
certificate for \eqref{eq:j-t}--\eqref{eq:split} and
\eqref{eq:disc}, an exact cyclotomic-orbit and ideal-rank certificate for the
chordal transverse divisor, an independent enumeration of the relevant
subgroups modulo six, their exact expected outputs, and a SHA-256 manifest.
The paper's
\texttt{make check} target replays all three checks, verifies the manifest, lints
the source, builds the PDF, and rejects TeX warnings.

\scriptsize
\begin{multicols}{2}
\raggedcolumns
\sloppy
\raggedright
\begin{thebibliography}{99}

\bibitem{ACT}
D.~Allcock, J.~A.~Carlson, and D.~Toledo,
\emph{The moduli space of cubic threefolds as a ball quotient},
Mem. Amer. Math. Soc. 209 (2011), no.~985;
doi:10.1090/S0065-9266-10-00591-0.

\bibitem{RebolledoWuthrich}
M.~Rebolledo and C.~Wuthrich,
\emph{A moduli interpretation for the non-split Cartan modular curve},
Glasg. Math. J. 60 (2018), no.~2, 411--434.

\bibitem{CMGHL}
S.~Casalaina-Martin, S.~Grushevsky, K.~Hulek, and R.~Laza,
\emph{Complete moduli of cubic threefolds and their intermediate Jacobians},
Proc. Lond. Math. Soc. 122 (2021), 259--316.

\bibitem{Hartlieb}
M.~Hartlieb,
\emph{Special subvarieties in the locus of intermediate Jacobians of cubic
threefolds}, Math. Z. 310 (2025), article 52.

\bibitem{LZ}
E.~Looijenga and Y.~Zi,
\emph{Monodromy and period map of the Winger pencil},
Nagoya Math. J. 251 (2023), 561--602.

\bibitem{Roulleau}
X.~Roulleau,
\emph{The Fano surface of the Klein cubic threefold},
J. Math. Kyoto Univ. 49 (2009), no.~1, 113--129.

\bibitem{Paulhus}
J.~Paulhus,
\emph{Elliptic factors in Jacobians of hyperelliptic curves with certain
automorphism groups},
Open Book Ser. 1 (2013), 487--505;
doi:10.2140/obs.2013.1.487.

\bibitem{RuddEpilogue}
T.~Rudd,
\emph{Irrationality of cubic threefolds after one stabilization},
preprint, 2026.

\bibitem{VGY}
B.~van Geemen and T.~Yamauchi,
\emph{On intermediate Jacobians of cubic threefolds admitting an
automorphism of order five},
Pure Appl. Math. Q. 12 (2016), no.~1, 141--164.

\end{thebibliography}
\end{multicols}

\end{document}
